\documentclass{article}
\usepackage[utf8]{inputenc}
\usepackage[default]{raleway}
\usepackage{titlesec, comment, tabularx, makecell, listings, array, setspace, geometry, graphicx, xcolor, xparse, fancyvrb, relsize, fancyhdr, booktabs, hyperref}
%\geometry{a4paper, left=2cm, right=2cm, top=2cm, bottom=2.5cm}
\renewcommand{\headrulewidth}{0pt}

% ----------------------------- Definizione tabella ---------------------------

\newcolumntype{C}[1]{>{\centering\arraybackslash}m{#1}}

% ------------------------------Metadati indice --------------------------------
\title{\textbf{\fontsize{30}{6}\selectfont Indice}}
\author{\fontsize{14}{6}\selectfont ByteOps}
\date{\today}

% -----------------------------Creazione footer --------------------------------
\pagestyle{fancy}
\fancyhf{}
\renewcommand{\footrulewidth}{0.4pt}
\lfoot{
    \parbox[c]{2cm}{\includegraphics[width=2cm]{../../Images/logo.png}}
    \textcolor[RGB]{120, 120, 120}{$\cdot$ Analisi dei rischi}
}
\rfoot{\thepage}

% --------------------------Modifica formato hyperlinks ------------------------
\hypersetup{
    colorlinks=true,
    linkcolor=black,
    filecolor=black,      
    pdftitle={Analisi dei rischi},
    pdfpagemode=FullScreen,
}



\begin{document}
\pagestyle{fancy}
\begin{center}
\includegraphics[width = 0.7\textwidth]{../../Images/logo.png} \\
\vspace{0.2cm}
\textcolor[RGB]{60, 60, 60}{\textit{ByteOps.swe@gmail.com}} \\
\vspace{1cm}
\fontsize{16}{6}\selectfont Analisi dei rischi \\
\vspace{0.5cm}
\end{center}

\section*{Informazioni documento}
\def\arraystretch{1.2}
\begin{tabular}{>{\raggedleft\arraybackslash}p{0.2\textwidth}|>{\raggedright\arraybackslash}p{0.6\textwidth}c}
\hline
\addlinespace
    \textbf{Redattori} & A. Barutta \\ & R. Smanio \\ & N. Preto \vspace{10pt} \\
    \textbf{Verificatori} & E. Hysa \\ & L. Skenderi \\ & D. Diotto \vspace{10pt} \\
    \textbf{Amministratore} & F. Pozza \vspace{10pt} \\
    \textbf{Destinatari} & T. Vardanega \\ & R. Cardin \vspace{10pt} \\
    \textbf{Partecipanti} & A. Barutta \\ & E. Hysa \\ & R. Smanio \\ & D. Diotto \\ & F. Pozza \\ & L. Skenderi \\ & N. Preto \vspace{10pt} \\
\end{tabular}
\pagebreak 

% ------------------------- Changelog ----------------------------

\begin{tabular}{|C{2.5cm}|C{2.5cm}|C{2.5cm}|C{2.5cm}|C{2.5cm}|}
    \hline
    \textbf{Versione} & \textbf{Data} & \textbf{Autore} & \textbf{Verificatore} & \textbf{Dettaglio} \\
    \hline \hline
    0.0.1 & 03/11/2023 & Nome autore & Nome verificatore & Scrittura delle sezioni 1 e 2 \\
    \hline

    
\end{tabular}

\pagebreak

% ------------------------- Generazione automatica indice ----------------------
\setstretch{1.5}
\maketitle
\thispagestyle{fancy}
\tableofcontents
\setstretch{1.2}
\pagebreak

% ---------------------------- Inizio documento -------------------------------


\section{Obiettivo}
Il seguente documento è stato redatto con l'obiettivo di valutare e disaminare i potenziali rischi che il gruppo potrebbe affrontare nel corso delle attività di sviluppo
del capitolato. In questo documento veranno esposti i vari rischi identificati, accompagnati da una relativa previsione di impatto che ciascuno di essi potrebbe avere e una possibile soluzione.


\section{Rischi}
\subsection{Rallentamento delle attività}
Uno dei principali rischi che potrebbero manifestarsi è l'armonizzazione tra le attività personali con quelle attività progettuali. Tale rischio
potrebbe causare un rallentamento nell'esecuzione dei compiti assegnati ai membri del gruppo, con conseguente rallentamento del processo di sviluppo. 
Si prevede che questo rischio si intensificherà durante il periodo di sessione invernale 2023-2024, a causa degli impegni legati agli esami.\\ 
Per quanto riguarda le possibili soluzioni per attenuare il problema, si propone: 
\begin{itemize}
    \item Una buona organizzazione tra i membri del gruppo per garantire una distribuzione equa del lavoro da svolgere.
    \item La creazione di un ambiente di lavoro asincrono, in modo da permettere lo svolgimento dei compiti assegnati da parte di ciascun membro del gruppo secondo le proprie tempistiche.
\end{itemize}

\subsection{Apprendimento ed utilizzo delle nuove tecnologie}
L’apprendimento e l'implementazione delle tecnologie proposte possono rappresentare un rischio considerevole per lo 
sviluppo di un progetto, in quanto esista la possibilità che lo studio accurato di queste tecnologie richieda più tempo del previsto. \\
Per quanto riguarda le possibili soluzioni per attenuare il problema, si propone:
\begin{itemize}
    \item Pianificazione accurata del tempo necessario per l’assimilazione e l’implementazione \\delle nuove tecnologie.
    \item Sfruttare le competenze individuali di ciascun membro del gruppo, al fine di accelerare \\ il processo di apprendimento.
    \item Approfittare, ove possibile, delle opportunità di formazione offerte dall’azienda proponente, in modo da risolvere eventuali dubbi e consolidare i concetti appresi.
\end{itemize}








\end{document} 